\documentclass[11pt,a4paper]{article}
\usepackage[T1]{fontenc}
\usepackage[utf8]{inputenc}
\usepackage{lmodern}
\usepackage[margin=27mm,headheight=15pt]{geometry}
\usepackage{amsmath,amssymb,amsthm,mathtools}
\usepackage{booktabs,tabularx,array,longtable}
\usepackage{microtype}
\usepackage{xcolor}
\usepackage{enumitem}
\usepackage{fancyhdr}
\usepackage{hyperref}
\definecolor{inkblue}{RGB}{28,57,83}
\hypersetup{colorlinks=true,linkcolor=inkblue,citecolor=inkblue,urlcolor=inkblue,
 pdftitle={Three Consecutive Square-Pyramidal Numbers: Exact Frobenius Formulas, Gap Laws, and Algebraic Inversion},
 pdfauthor={Research report prepared for Vladimir Reshetnikov}}
\pagestyle{fancy}
\fancyhf{}
\fancyhead[L]{\small\textsc{An explicit solution for OEIS A069762}}
\fancyhead[R]{\small Research report}
\fancyfoot[C]{\thepage}
\renewcommand{\headrulewidth}{0.3pt}
\setlength{\parskip}{3pt}
\setlength{\emergencystretch}{2em}
\numberwithin{equation}{section}
\newtheorem{theorem}{Theorem}[section]
\newtheorem{lemma}[theorem]{Lemma}
\newtheorem{proposition}[theorem]{Proposition}
\newtheorem{corollary}[theorem]{Corollary}
\theoremstyle{definition}
\newtheorem{definition}[theorem]{Definition}
\newtheorem{remark}[theorem]{Remark}
\newcommand{\N}{\mathbb N_0}
\newcommand{\Z}{\mathbb Z}
\newcommand{\Q}{\mathbb Q}
\newcommand{\R}{\mathbb R}
\newcommand{\Ap}{\operatorname{Ap}}
\newcommand{\Gap}{\operatorname{Gap}}
\newcommand{\eps}{\varepsilon}
\newcommand{\ind}{\mathbf 1}
\newcommand{\E}{\mathbb E}
\newcommand{\Prob}{\mathbb P}
\newcommand{\dd}{\,\mathrm d}
\newcommand{\ceil}[1]{\left\lceil#1\right\rceil}
\newcommand{\floor}[1]{\left\lfloor#1\right\rfloor}
\newcommand{\Bern}{\mathrm B}
\title{\vspace{-1.0cm}\color{inkblue}\LARGE\bfseries
Three Consecutive Square-Pyramidal Numbers\\[5pt]
\Large Exact Frobenius Formulas, Gap Laws,\\and Algebraic Inversion}
\author{\normalsize Research report prepared for Vladimir Reshetnikov}
\date{October 1, 2026}

\begin{document}
\maketitle
\thispagestyle{empty}
\begin{abstract}
Let $s_n=n(n+1)(2n+1)/6$ and let $F_n$ be the largest integer not in
$\langle s_n,s_{n+1},s_{n+2}\rangle$. This is OEIS A069762. We give an explicit
solution: $F_n$ is a degree-five quasipolynomial of least eventual period six,
with exactly six exceptional indices, the last being $23$. In particular,
$F_n=n^5/9+8n^4/9+O(n^3)$, and all subsequent terms are given exactly.

The proof does considerably more than evaluate a maximum. Two elementary gcd
reductions expose a linear-size Ap\'ery set controlled by the relation
$kC-\ell A=5B$, with $k,\ell\leq6$. This yields a direct membership and
representation algorithm, exact genus and symmetry formulas, and exact
quasipolynomials for every fixed power sum of the gaps. A general transfer
principle converts limits of normalized Ap\'ery sets into limits of gap
measures. Applied here, a uniformly chosen gap divided by $F_n$ converges to
the density $2(1-x)$ on $[0,1]$, with a uniform $O(n^{-1})$ error in the
cumulative distribution. The ordinary generating function has minimal
denominator $(1-x)^6(1+x)^4(1+x+x^2)^4$, hence a minimal eventual recurrence
of order $18$.

Finally, each residue-class polynomial admits a convergent all-orders
Puiseux inverse. We give its first coefficients, a coefficient-extraction
formula, a uniform sufficient convergence region, and an exact rule for
recovering the discrete inverse. All main claims are proved in the text;
exact computational checks and symbolic moment-generation programs accompany
the report. The explicit formulas are proposed as additions to the inspected
OEIS entry; literature priority is not certified.
\end{abstract}
\vfill
\noindent\small\textbf{Scope of the claim.} This is an independently derived research
report, not a claim that a famous named conjecture has been settled. General
eventual quasipolynomiality is established literature. The contribution is the
explicit formula, complete normal form, quantitative gap law, and effective
consequences for this particular canonical family. No Lean or Rocq
formalization of these results is claimed.
\clearpage
\tableofcontents
\clearpage

\section{The target, the literature, and the proof boundary}

The inspected version of OEIS A069762 defines the Frobenius numbers generated
by three consecutive square-pyramidal numbers, supplies values for
$n=2,\ldots,35$, and gives the example
$F(5,14,30)=51$. It contains no explicit general formula or asymptotic
formula \cite{oeis}. Its displayed definition is not labelled a conjecture.
Thus the precise problem addressed here is to supply a missing explicit
solution and its proof, rather than to reclassify an OEIS omission as a
well-established open conjecture.

This is a natural parametric Frobenius problem, not an artificial sequence
chosen solely because a few terms interpolate. Frobenius numbers for three
consecutive squares and cubes were explicitly analyzed by Lepilov,
O'Rourke, and Swanson \cite{los}. Related work of Robles-P\'erez and Rosales
studies consecutive triangular and tetrahedral numbers \cite{rr}. Those
families provide relevant context, but square-pyramidal numbers are a
different polynomial family. Roune and Woods proved eventual
quasipolynomiality for parametric Frobenius problems with at most three
generators \cite{rw}; Shen established the broader polynomial-parametric
result \cite{shen}. We do not claim eventual quasipolynomiality itself as new.

The source check for this report used the OEIS entry, the cited primary
literature, searches for A069762 and for Frobenius numbers of square-pyramidal
numbers, and a scoped inspection of the ProveIt repository. No publication
containing the formula below was located. This is a bounded literature
check, not a proof of priority. A prior equivalent formula, perhaps under
different terminology, could change the novelty assessment without changing
the mathematical statements proved here.

The connection to ProveIt is methodological rather than a dependency on an
unverified repository theorem. Its README separates exploratory computations
from the trusted boundary of formal mathematics \cite{proveit}. We follow
that separation: computation suggested the formula, but the proof is the
explicit normal form developed below. The accompanying programs provide
independent finite checks and algebraic audits. They are neither a proof of
literature novelty nor a kernel-checked formalization. The repository snapshot
used for this inspection is recorded in the bibliography; the repository
was not modified.

\paragraph{What is established here.}
The main output is an exact answer for \emph{every} $n\geq2$, not just a fitted
asymptotic equivalent. The proof also determines all missing integers through
a constant-description family of residue representatives. The distribution
and moment results explain the global shape of the gaps, while the inverse
results address the associated threshold problem. These additional results
make the normal form reusable beyond a single OEIS formula.

\section{An exact formula for A069762}\label{sec:main}

For $n\geq2$, put
\begin{equation}\label{eq:defs}
s_n=\sum_{j=1}^n j^2=\frac{n(n+1)(2n+1)}6,
\qquad S_n=\langle s_n,s_{n+1},s_{n+2}\rangle,
\qquad F_n=\max(\N\setminus S_n).
\end{equation}
Here $\langle a,b,c\rangle=\{ua+vb+wc:u,v,w\in\N\}$.
Define the periodic integer
\begin{equation}\label{eq:h}
h_n=3\ind_{\{n\text{ odd}\}}+2\ind_{\{n\equiv2\pmod3\}}.
\end{equation}
Its values for residues $0,1,2,3,4,5$ modulo six are
$0,3,2,3,0,5$.
For $h\in\{0,2,3,5\}$ define
\begin{align}\label{eq:Ph}
P_h(X)=\frac1{36}\bigl(&4X^5+32X^4+(101-10h)X^3\notag\\
 &+(175-45h)X^2+(102-65h)X-36-30h\bigr).
\end{align}
Equivalently,
\begin{equation}\label{eq:compactP}
P_h(X)=P_0(X)-\frac{5h}{36}(X+1)(X+2)(2X+3).
\end{equation}

\begin{samepage}
\begin{theorem}[Exact Frobenius formula]\label{thm:main}
For every integer $n\geq2$,
\begin{equation}\label{eq:F}
\boxed{\ F_n=P_{h_n}(n)+\eps_n\ },
\end{equation}
where $\eps_n=0$ except at the six indices in Table~\ref{tab:exceptions}.
Consequently $F_n=P_{h_n}(n)$ for all $n\geq24$, and $24$ is the least
uniform cutoff for this quasipolynomial. Its least eventual period is six.
\end{theorem}
\end{samepage}

\begin{table}[ht]
\centering
\caption{The four polynomial branches. The common numerator begins
$4n^5+32n^4$ and is divided by $36$.}\label{tab:branches}
\begin{tabular}{ccrrrr}
\toprule
$n\bmod6$ & $h_n$ & coefficient of $n^3$ & of $n^2$ & of $n$ & constant\\
\midrule
$0,4$ & 0 & 101 & 175 & 102 & $-36$\\
$1,3$ & 3 & 71 & 40 & $-93$ & $-126$\\
$2$   & 2 & 81 & 85 & $-28$ & $-96$\\
$5$   & 5 & 51 & $-50$ & $-223$ & $-186$\\
\bottomrule
\end{tabular}
\end{table}

\begin{table}[ht]
\centering
\caption{All exceptions: the correction is positive in each case.}\label{tab:exceptions}
\begin{tabular}{rrrr}
\toprule
$n$ & $P_{h_n}(n)$ & $\eps_n$ & $F_n$\\
\midrule
2 & 41 & 10 & 51\\
3 & 151 & 40 & 191\\
5 & 1009 & 315 & 1324\\
11 & 32553 & 1612 & 34165\\
17 & 238451 & 3135 & 241586\\
23 & 980251 & 2400 & 982651\\
\bottomrule
\end{tabular}
\end{table}

The exact formula immediately yields
\begin{equation}\label{eq:asympF}
F_n=\frac{n^5}{9}+\frac{8n^4}{9}
 +\frac{101-10h_n}{36}n^3
 +\frac{175-45h_n}{36}n^2
 +\frac{102-65h_n}{36}n-1-\frac{5h_n}{6}
\end{equation}
for $n\geq24$. This is a complete finite asymptotic expansion with periodic
coefficients: there is no remainder after the constant term. On each residue
class it is simply a polynomial identity. In particular,
$F_n\sim n^5/9$. A single constant coefficient at order $n^3$ would be
incorrect; the first oscillatory term occurs at precisely that order.

The exact formula and the exception list are proved in
Sections~\ref{sec:reduction}--\ref{sec:max}. Minimality of the period is
proved in Section~\ref{sec:gf}.

\section{Ap\'ery sets and two elementary reductions}\label{sec:reduction}

\subsection{The residue description of all gaps}

A numerical semigroup is an additive submonoid of $\N$ with finite
complement. Positive integers with gcd one generate such a semigroup:
modulo any one generator, their residues generate the entire finite cyclic
group, and a submonoid of a finite group is a subgroup. Thus each residue
has a nonnegative representative; adding multiples of the chosen generator
proves cofiniteness.

For a numerical semigroup $S$ and $m\in S\setminus\{0\}$, let $w_r$ be the
least element of $S$ congruent to $r$ modulo $m$, for $0\leq r<m$. The
Ap\'ery set is
\[
\Ap(S;m)=\{w_0,\ldots,w_{m-1}\}.
\]
In residue $r$, the gaps are exactly
\begin{equation}\label{eq:resgaps}
r,\ r+m,\ldots,\ w_r-m,
\end{equation}
with the list empty when $w_r=r$. In particular $w_0=0$, so zero is never
counted as a gap. It follows directly that
\begin{equation}\label{eq:apfacts}
F(S)=\max_r w_r-m,
\qquad
 g(S)=\#(\N\setminus S)=\frac1m\sum_r w_r-\frac{m-1}{2},
\end{equation}
and, as a formal power series,
\begin{equation}\label{eq:hilbert}
H_S(z):=\sum_{s\in S}z^s
=\frac{\sum_{w\in\Ap(S;m)}z^w}{1-z^m}.
\end{equation}
These facts are standard in numerical-semigroup theory; the short derivation
above fixes all indexing conventions used in this report. The gcd reduction
below is the mechanism behind the classical Johnson Frobenius reduction
recalled in \cite{rr}, but we prove the full normal form needed here.

\subsection{A rectangular lifting lemma}

\begin{lemma}[Two-stage normal form]\label{lem:lift}
Let $a,b,c$ be positive integers with gcd one. Set
\[
d=\gcd(a,b),\quad e=\gcd(b,c),\quad
 A=a/d,\quad B=b/(de),\quad C=c/e.
\]
Then $d,e$ are coprime, the displayed quantities are positive integers, and
with $T=\langle A,B,C\rangle$ one has the disjoint decomposition
\begin{equation}\label{eq:rect}
\langle a,b,c\rangle
=\bigsqcup_{\substack{0\leq i<e\\0\leq j<d}}
\bigl(ia+jc+de\,T\bigr).
\end{equation}
The offsets give all residues modulo $de$. Moreover,
\begin{equation}\label{eq:liftap}
\Ap(\langle a,b,c\rangle;b)
=\{de\,w+ia+jc:w\in\Ap(T;B),\ 0\leq i<e,\ 0\leq j<d\}.
\end{equation}
Consequently,
\begin{align}
F(\langle a,b,c\rangle)
 &=deF(T)+(e-1)a+(d-1)c,\label{eq:johnson}\\
g(\langle a,b,c\rangle)
 &=de\,g(T)+\frac{(e-1)a+(d-1)c-de+1}{2}.\label{eq:genuslift}
\end{align}
\end{lemma}

\begin{proof}
Any common divisor of $d,e$ divides $a,b,c$, hence $\gcd(d,e)=1$ and
$de\mid b$. Also $\gcd(d,c)=\gcd(e,A)=1$.
In a representation $ua+vb+wc$, reduce the coefficient $w$ modulo $d$.
This gives uniquely a residue $0\leq j<d$ and a remainder in
$d\langle A,eB,c\rangle$. In that remainder reduce the coefficient of $A$
modulo $e$. This gives $0\leq i<e$ and a remainder in
$e\langle A,B,C\rangle$. Conversely, every expression on the right of
\eqref{eq:rect} is a nonnegative combination of $a,b,c$, since
$deA=ea$, $deB=b$, and $deC=dc$.

If two offsets are equal modulo $de$, reduction modulo $d$ first forces
their $j$ values to agree, because $c$ is invertible modulo $d$. After
division by $d$, reduction modulo $e$ forces their $i$ values to agree,
because $A$ is invertible modulo $e$. This proves disjointness and shows
that all $de$ residues occur. For each fixed offset, minimizing further
modulo $b=deB$ is exactly the Ap\'ery minimization in $T$ modulo $B$.
This proves \eqref{eq:liftap}.

Taking the maximum in \eqref{eq:liftap} proves \eqref{eq:johnson}.
Its sum is
\[
(de)^2\sum_{w\in\Ap(T;B)}w
+Bde\left(\frac{e-1}{2}a+\frac{d-1}{2}c\right).
\]
Apply \eqref{eq:apfacts}, with modulus $deB$, to obtain
\eqref{eq:genuslift}.
\end{proof}

\subsection{The gcds of consecutive square-pyramidal numbers}

Throughout the remainder of the paper, unless otherwise specified,
\begin{equation}\label{eq:params}
\begin{gathered}
a=s_n,\qquad b=s_{n+1},\qquad c=s_{n+2},\\
k=\gcd(n+1,6),\quad \ell=\gcd(n+2,6),\qquad
 d=\frac{n+1}{k},\quad e=\frac{n+2}{\ell},\\
A=\frac{k n(2n+1)}6,\qquad
B=\frac{k\ell(2n+3)}6,\qquad
C=\frac{\ell(n+3)(2n+5)}6.
\end{gathered}
\end{equation}
The relation between these symbols and Lemma~\ref{lem:lift} is exact.

\begin{lemma}\label{lem:gcd}
One has $\gcd(a,b)=d$, $\gcd(b,c)=e$, and $\gcd(a,b,c)=1$.
The values of $k,\ell,B$ are given by Table~\ref{tab:params}. In particular,
$\gcd(k,\ell)=1$, $B$ is odd, and $B\equiv1\pmod{k}$.
\end{lemma}

\begin{proof}
For $m=n+1$,
\[
\gcd(s_{m-1},s_m)=\gcd\left(\frac{(m-1)m(2m-1)}6,m^2\right).
\]
If a prime $p\geq5$ divides $m$, neither $m-1$ nor $2m-1$ is divisible by
$p$, so the $p$-adic valuation of this gcd is $v_p(m)$. If $2\mid m$,
the other two factors are odd and the valuation is $v_2(m)-1$; if
$3\mid m$, the valuation is $v_3(m)-1$. Primes not dividing $m$ do not
contribute. Thus the gcd is $m/\gcd(m,6)$, proving the two adjacent-gcd
formulas. A common divisor of $a,b,c$ divides both
$b-a=(n+1)^2$ and $c-b=(n+2)^2$, which are coprime. The table is obtained
by evaluating two gcds modulo six. Its remaining assertions follow at once.
\end{proof}

\begin{table}[ht]
\centering
\caption{The reduced parameters. The linear modulus $B$ replaces the
original cubic modulus $b$.}\label{tab:params}
\begin{tabular}{ccccc}
\toprule
$n\bmod6$ & $k$ & $\ell$ & $B$ & gaps of $\langle k,\ell\rangle$\\
\midrule
0 & 1 & 2 & $(2n+3)/3$ & none\\
1 & 2 & 3 & $2n+3$ & $\{1\}$\\
2 & 3 & 2 & $2n+3$ & $\{1\}$\\
3 & 2 & 1 & $(2n+3)/3$ & none\\
4 & 1 & 6 & $2n+3$ & none\\
5 & 6 & 1 & $2n+3$ & none\\
\bottomrule
\end{tabular}
\end{table}

The decisive identities are
\begin{equation}\label{eq:key}
\boxed{\ kC-\ell A=5B\ },\qquad
2A\equiv k\pmod B,\qquad 2C\equiv\ell\pmod B.
\end{equation}
For completeness, exact forms of the congruences are
\begin{equation}\label{eq:exactcong}
2A-k=B\frac{2(n-1)}{\ell},\qquad
2C-\ell=B\frac{2(n+4)}k.
\end{equation}
Both quotients are integers in each row of Table~\ref{tab:params}.
The first identity follows by subtracting the two quadratic expressions in
\eqref{eq:params}. No approximation is involved in these reductions.

\section{A bounded-coefficient Ap\'ery theorem}\label{sec:apery}

\subsection{A reusable reduction principle}

\begin{lemma}[Least auxiliary residue principle]\label{lem:aux}
Let $A,B,C,k,\ell,H$ be positive integers with $\gcd(k,\ell)=1$. Suppose
\[
kC-\ell A=HB,
\qquad tA\equiv k\pmod B,\qquad tC\equiv\ell\pmod B,
\qquad \gcd(t,B)=1.
\]
Assume $A>H(k-1)$. For each residue modulo $B$, take the least nonnegative
member $L$ of $\langle k,\ell\rangle$ in that residue. Let
\[
0\leq v(L)<k,\quad \ell v(L)\equiv L\pmod k,
\qquad u(L)=\frac{L-\ell v(L)}k.
\]
Then
\begin{equation}\label{eq:weightgeneral}
\Ap(\langle A,B,C\rangle;B)
=\left\{Au(L)+Cv(L)\right\}
=\left\{\frac{AL+HBv(L)}k\right\},
\end{equation}
where $L$ runs through the selected $B$ auxiliary representatives.
\end{lemma}

\begin{proof}
Every candidate for a least residue representative can be written using
$A,C$ only, because a positive coefficient of $B$ can be removed. Replacing
$kC$ by $\ell A$ lowers the value by $HB$ without changing its residue
modulo $B$. Thus a minimizing representation has its coefficient of $C$
in $[0,k-1]$.

For a normalized representation $uA+vC$, set $L=ku+\ell v$.
The congruence class of $L$ modulo $B$ is $t(uA+vC)$, so auxiliary residues
and value residues correspond bijectively. For any representable $L$, the
normalized coefficient $v(L)$ is unique, and its coefficient $u(L)$ is
nonnegative: normalize any nonnegative representation of $L$ in
$\langle k,\ell\rangle$ by reducing its coefficient of $\ell$ modulo $k$.
The relation gives the second formula in \eqref{eq:weightgeneral}.

Let $L$ be the least representable auxiliary integer in a fixed residue.
Every other representable auxiliary integer there is $L+qB$, with $q\geq1$.
The corresponding difference of normalized weights is
\[
\frac{B}{k}\left(Aq+H\bigl(v(L+qB)-v(L)\bigr)\right)
\geq\frac{B}{k}\bigl(A-H(k-1)\bigr)>0.
\]
Hence the least auxiliary integer gives the least value. There is one such
integer in every residue because $\langle k,\ell\rangle$ has gcd one.
This proves both completeness and minimality.
\end{proof}

The point of this lemma is not merely to reduce the number of generators.
It replaces an unbounded coefficient search by a bounded coefficient
$0\leq v<k$ and a least-residue problem in a fixed, very small semigroup.
The constant $H$ need not equal five, and the parameter family need not be
square-pyramidal.

\subsection{The complete reduced Ap\'ery set}

For the parameters \eqref{eq:params}, define
\begin{equation}\label{eq:Lset}
\mathcal L_n=
\begin{cases}
\bigl(\{0,1,\ldots,B-1\}\setminus\{1\}\bigr)\cup\{B+1\},
   & n\equiv1,2\pmod6,\\
\{0,1,\ldots,B-1\},&\text{otherwise}.
\end{cases}
\end{equation}
For $L\in\mathcal L_n$, put
\begin{equation}\label{eq:vuw}
\begin{gathered}
v(L)\in\{0,\ldots,k-1\},\qquad \ell v(L)\equiv L\pmod k,\\
u(L)=\frac{L-\ell v(L)}k,\qquad
w(L)=Au(L)+Cv(L)=\frac{AL+5Bv(L)}k.
\end{gathered}
\end{equation}
When $k=1$, the definition means simply $v(L)=0$.

\begin{theorem}[Explicit Ap\'ery normal form]\label{thm:apery}
For every $n\geq2$,
\begin{equation}\label{eq:smallap}
\Ap(\langle A,B,C\rangle;B)=\{w(L):L\in\mathcal L_n\},
\end{equation}
and
\begin{equation}\label{eq:fullap}
\boxed{\ \Ap(S_n;b)=
\{de\,w(L)+ia+jc:L\in\mathcal L_n,\ 0\leq i<e,\ 0\leq j<d\}\ }.
\end{equation}
All the representatives in \eqref{eq:fullap} have different residues
modulo $b$. Thus this formula gives every gap by \eqref{eq:resgaps}.
\end{theorem}

\begin{proof}
Table~\ref{tab:params} shows that the only missing auxiliary integer is
$1$, precisely in the two rows with generators $2,3$. This proves that
\eqref{eq:Lset} consists of the least representable auxiliary integers in
each residue modulo $B$.

For $n\geq3$, one has $A>5(k-1)$. For $k=1$ the assertion is immediate.
For $k=2$ the smallest relevant $n$ is $3$, giving $A=7>5$; for $k=3$
it is $8$, giving $A=68>10$; and for $k=6$ it is $5$, giving
$A=55>25$. In each of these rows $A$ increases with $n$.
Apply Lemma~\ref{lem:aux} with $H=5$ and $t=2$, using that $B$ is odd.
This proves \eqref{eq:smallap} for $n\geq3$.

At $n=2$, the parameters are $(A,B,C)=(5,7,15)$ and
$\mathcal L_2=\{0,2,3,4,5,6,8\}$. The corresponding weights are
$0,15,5,30,20,10,25$, exactly the set
$\{0,5,10,15,20,25,30\}=\Ap(\langle5,7\rangle;7)$.
Since $15=3\cdot5$, this is also the required Ap\'ery set for
$\langle5,7,15\rangle$. This explicit check covers the only value for
which the sufficient inequality in Lemma~\ref{lem:aux} fails.
Finally apply Lemma~\ref{lem:lift}.
\end{proof}

\section{Taking the maximum and explaining all six exceptions}\label{sec:max}

Write $M=\max_{L\in\mathcal L_n}w(L)$. The Frobenius number of the reduced
semigroup is $M-B$, and \eqref{eq:johnson} gives
\begin{equation}\label{eq:FfromM}
F_n=de(M-B)+(e-1)a+(d-1)c.
\end{equation}
The only task remaining for the largest gap is to identify $M$.

\begin{table}[ht]
\centering
\caption{Exact maximum candidates; both candidates are needed at the small
exceptional indices.}\label{tab:max}
\begin{tabular}{cp{0.71\textwidth}}
\toprule
$n\bmod6$ & $M$\\
\midrule
$0,4$ & $A(B-1)$\\[4pt]
$1$ & $\max\{A(B+1)/2,\ A(B-5)/2+C\}$\\[4pt]
$2$ & $\max\{A(B-1)/3+C,\ A(B-7)/3+2C\}$\\[4pt]
$3$ & $\max\{A(B-1)/2,\ A(B-3)/2+C\}$\\[4pt]
$5$ & $\max\{A(B-1)/6,\ A(B-7)/6+5C\}$\\
\bottomrule
\end{tabular}
\end{table}

To verify Table~\ref{tab:max}, fix $v$. The weight $Au+Cv$ increases with
$u$, so only the largest allowed $u$ matters. Since $B\equiv1\pmod k$,
that largest value, before the possible extra auxiliary integer $B+1$, is
\begin{equation}\label{eq:U}
U_v=\frac{B-1}{k}-\ceil{\frac{\ell v}{k}}.
\end{equation}
For residue $1$, the extra $B+1$ has $v=0$ and supplies the first candidate.
For residue $2$, it has $v=1$ and supplies the first candidate; the $v=0$
candidate is smaller by $C$. For residue $5$, all $v=1,\ldots,5$ have
$u=(B-7)/6$, so only $v=5$ competes with $v=0$. The other rows are
immediate. This derives every candidate rather than inferring it from data.

The second candidate minus the first is, respectively,
\begin{equation}\label{eq:signs}
\begin{array}{c|c|c}
n\bmod6 & \text{difference} & \text{first candidate wins from}\\ \hline
1 & C-3A=(5B-3A)/2 & n=7\\
2 & C-2A=(5B-4A)/3 & n=8\\
3 & C-A=(5B-A)/2 & n=9\\
5 & 5C-A=(25B-A)/6 & n=29.
\end{array}
\end{equation}
Here ``from'' means along the indicated residue class. The signs reduce to
quadratic inequalities. In the first row the numerator $5B-3A$ is
$-2n^2+9n+15$; in the second, $5B-4A=-4n^2+8n+15$; in the third,
$5B-A=(-2n^2+9n+15)/3$; and in the fourth,
$25B-A=-2n^2+49n+75$. These are negative at and beyond the listed
thresholds. They are positive at the earlier allowed indices, namely
$2$, $3$, and $5,11,17,23$ in the last three rows. There are no earlier
allowed indices in the first row because $n\geq2$.

Substituting the first candidates into \eqref{eq:FfromM} and simplifying
gives the four polynomials in Table~\ref{tab:branches}. At an exceptional
index the correction is $de$ times the positive difference in
\eqref{eq:signs}. This gives exactly Table~\ref{tab:exceptions}, proving
the formula and the sharp cutoff in Theorem~\ref{thm:main}.

\paragraph{Why early interpolation is hazardous.}
The exceptional behavior is caused by a genuine change of maximizing
Ap\'ery element, not by an indexing error. In residue class five the
exceptional branch persists through $n=23$. A polynomial fitted only to
small values in that residue can therefore look consistent while describing
the wrong eventual branch. The normal form identifies the exact competition
and proves where it ends.

\section{All gaps, constructive membership, and a worked example}\label{sec:membership}

\subsection{Membership without constructing the Ap\'ery set}

The normal form yields a bounded-operation membership procedure. Given an
integer $x$, choose
\begin{equation}\label{eq:decode}
\begin{gathered}
0\leq j<d,\qquad jc\equiv x\pmod d,\qquad
 y=\frac{x-jc}{d},\\
0\leq i<e,\qquad iA\equiv y\pmod e,\qquad
 z=\frac{y-iA}{e}.
\end{gathered}
\end{equation}
The congruences have unique solutions because of the coprimalities proved
in Lemma~\ref{lem:lift}. A congruence modulo one simply selects zero.
Let $L_0$ be the least residue of $2z$ modulo $B$. Set $L=L_0$, except that
$L=B+1$ when $L_0=1$ and $n\equiv1,2\pmod6$.

\begin{corollary}[Membership and representation certificate]\label{cor:membership}
With the definitions above,
\begin{equation}\label{eq:membership}
x\in S_n\quad\Longleftrightarrow\quad z\geq w(L).
\end{equation}
When membership holds, $q=(z-w(L))/B$ is a nonnegative integer and
\begin{equation}\label{eq:repr}
\boxed{\ x=(i+e\,u(L))a+q b+(j+d\,v(L))c\ }.
\end{equation}
\end{corollary}

\begin{proof}
The decoding recovers the unique offset class in \eqref{eq:rect}, even when
$x$ or an intermediate quotient is negative. The congruences in
\eqref{eq:key} imply $2w(L)\equiv L\equiv2z\pmod B$, so $w(L)$ is the
least reduced representative in the residue of $z$. Since $2$ is invertible
modulo $B$, the difference $z-w(L)$ is a multiple of $B$. Theorem~\ref{thm:apery}
now proves \eqref{eq:membership}; substituting the definitions proves
\eqref{eq:repr}. All its coefficients are nonnegative.
\end{proof}

For fixed $n$, the two modular inverses can be precomputed. Each subsequent
membership query uses a bounded number of exact arithmetic operations,
independent of $b$ and $F_n$. This is an arithmetic-operation statement, not
a claim that arbitrarily large integers take constant bit time. The program
\texttt{representation} in \texttt{verify.py} implements the formula and
returns either a nonnegative coefficient triple or \texttt{None}.

\subsection{The complete generating function of the semigroup}

Writing $R_m(t)=1+t+\cdots+t^{m-1}$, \eqref{eq:fullap} gives
\begin{equation}\label{eq:fullhilbert}
H_{S_n}(z)=
\frac{R_e(z^a)R_d(z^c)}{1-z^b}
\sum_{L\in\mathcal L_n}z^{de\,w(L)}.
\end{equation}
Consequently the \emph{gap polynomial} is exactly
\begin{equation}\label{eq:gappoly}
\sum_{t\in\Gap(S_n)}z^t=\frac1{1-z}-H_{S_n}(z).
\end{equation}
The right side is a polynomial despite its displayed rational form: the
residue description proves that all sufficiently large coefficients cancel.
This is a full description of the complement, not just an estimate for its
largest element.

\subsection{The exceptional case \texorpdfstring{$n=5$}{n=5}}

Here $(a,b,c)=(55,91,140)$, $(d,e)=(1,7)$, $(k,\ell)=(6,1)$, and
$(A,B,C)=(55,13,20)$. The relation is $6\cdot20-55=5\cdot13$.
For $L=0,\ldots,12$, the reduced weights are
\[
0,20,40,60,80,100,55,75,95,115,135,155,110.
\]
Their maximum is $155$ at $L=11$, not the eventual-branch candidate
$110$ at $L=12$. Hence
\[
F(\langle55,13,20\rangle)=155-13=142,
\qquad F_5=7\cdot142+6\cdot55=1324.
\]
The largest original Ap\'ery representative is
$7\cdot155+6\cdot55=1415$, and subtracting $b=91$ gives the same
Frobenius number. The reduced sum is $1040$, so its genus is
$1040/13-6=74$; the original genus is
$7\cdot74+(330-7+1)/2=680$.
The entire exception is therefore explained by a single visible change of
maximum in a thirteen-element set.

\section{The genus, symmetry, and every gap moment}\label{sec:moments}

\subsection{A genus formula with no exceptions}

Let $g_n=\#\Gap(S_n)$, and define the six-periodic sequence
\begin{equation}\label{eq:J}
J_n=(0,9,8,9,0,5)_{n\bmod6}.
\end{equation}

\begin{theorem}[Exact genus]\label{thm:genus}
For every $n\geq2$,
\begin{equation}\label{eq:genus}
\boxed{\
g_n=\frac{n(n+1)(4n^3+28n^2+73n+102)
             -5J_n(n+1)(n+2)}{72}\ }.
\end{equation}
There are no exceptional indices. In particular,
\begin{equation}\label{eq:genusexp}
g_n=\frac{n^5}{18}+\frac{4n^4}{9}+\frac{101n^3}{72}
+\frac{175-5J_n}{72}n^2+\frac{102-15J_n}{72}n-\frac{5J_n}{36}.
\end{equation}
Thus the first three coefficients are independent of the residue class.
\end{theorem}

\begin{proof}
Put $\eta=1$ in residues $1,2$ and $\eta=0$ otherwise. Directly from
\eqref{eq:Lset},
\[
\sum_{L\in\mathcal L_n}L=\frac{B(B-1)}2+B\eta.
\]
The sum $V=\sum_{L\in\mathcal L_n}v(L)$ is
\begin{equation}\label{eq:V}
\begin{array}{c|rrrrrr}
n\bmod6&0&1&2&3&4&5\\ \hline
V&0&(B-3)/2&B-2&(B-1)/2&0&5(B-1)/2.
\end{array}
\end{equation}
For example, in residue $2$ the unmodified interval contributes $B-1$;
removing $L=1$ removes $v=2$, and adding $L=B+1$ adds $v=1$.
In residue $1$ the analogous change removes one. The other rows just sum
complete periods of $v(L)$.

Using \eqref{eq:vuw} and \eqref{eq:apfacts}, the reduced genus is
\begin{equation}\label{eq:redgenus}
g(T)=\frac Ak\left(\frac{B-1}{2}+\eta\right)
       +\frac5k V-\frac{B-1}{2}.
\end{equation}
Insert this into \eqref{eq:genuslift} and use
\eqref{eq:params}. The result in each of the six rows is
\eqref{eq:genus}. All identities remain valid at $n=2$ by the explicit
Ap\'ery set already proved in Theorem~\ref{thm:apery}.
\end{proof}

\begin{corollary}[Exact symmetry classification]\label{cor:symmetry}
The semigroup $S_n$ is symmetric if and only if
\begin{equation}\label{eq:symmetric}
n=2\quad\text{or}\quad n\equiv0,4\pmod6.
\end{equation}
For every other $n\geq2$, it is nonsymmetric, although
$g_n/F_n\longrightarrow1/2$.
\end{corollary}

\begin{proof}
A numerical semigroup is symmetric when for every integer $t$ exactly one
of $t,F-t$ belongs to it. Pairing the integers in $[0,F]$ shows the standard
equivalence with $2g=F+1$: both members of a pair cannot be in the
semigroup, and equality means there are no pairs with both members gaps.
Here direct subtraction gives
\begin{equation}\label{eq:defect}
2g_n-F_n-1=
\begin{cases}
0,&n\equiv0,4\pmod6,\\
\dfrac{5n(n+1)(n+2)}6-\eps_n,&n\equiv1,3\pmod6,\\
\dfrac{5(n+1)(n+2)(2n-1)}{18}-\eps_n,&n\equiv2\pmod6,\\
\dfrac{25(n+1)^2(n+2)}{18}-\eps_n,&n\equiv5\pmod6.
\end{cases}
\end{equation}
At the exceptional indices $2,3,5,11,17,23$, the actual defects are
$0,10,35,988,5415,17600$. Outside them the nonzero branches are positive.
This proves the classification. The asymptotic ratio follows from the
leading terms of \eqref{eq:asympF} and \eqref{eq:genusexp}.
\end{proof}

\subsection{A closed formula for every power sum}

For an integer $p\geq0$, define
\[
M_p(n)=\sum_{t\in\Gap(S_n)}t^p,
\qquad M_0(n)=g_n.
\]
Let $\Bern_j(X)$ be the Bernoulli polynomials in the convention
\[
\frac{u e^{Xu}}{e^u-1}=\sum_{j\geq0}\Bern_j(X)\frac{u^j}{j!},
\qquad \beta_j=\Bern_j(0),\quad \beta_1=-\frac12.
\]

\begin{proposition}[Ap\'ery--Bernoulli moment formula]\label{prop:bernoulli}
For a numerical semigroup $S$ and $m\in S\setminus\{0\}$,
\begin{equation}\label{eq:moment}
\boxed{\ \sum_{t\in\Gap(S)}t^p
=\frac{m^p}{p+1}\sum_{w\in\Ap(S;m)}
       \Bern_{p+1}(w/m)-\frac{\beta_{p+1}}{p+1}\ }.
\end{equation}
In particular, for $S=S_n$, the representatives in
Theorem~\ref{thm:apery} give a finite exact formula for every $M_p(n)$.
\end{proposition}

\begin{proof}
Write $w_r=r+m q_r$. Summing powers in the arithmetic progression
\eqref{eq:resgaps} gives
\[
\sum_{j=0}^{q_r-1}(r+mj)^p
=\frac{m^p}{p+1}
 \left(\Bern_{p+1}(w_r/m)-\Bern_{p+1}(r/m)\right).
\]
The Bernoulli multiplication identity
$\sum_{r=0}^{m-1}\Bern_{p+1}(r/m)=m^{-p}\beta_{p+1}$
follows immediately by summing the defining generating functions.
Now sum over $r$.
\end{proof}

The classical power-sum identity in this proposition is not a novelty claim.
What makes it effective here is the complete uniform description of its
Ap\'ery input. The next theorem gives a global consequence for all moment
orders, without requiring each polynomial to be separately guessed.

\begin{theorem}[Exact all-orders moment quasipolynomials]\label{thm:allmoments}
For every fixed integer $p\geq0$, $M_p(n)$ is a quasipolynomial of period
dividing six for all $n\geq2$. Each residue-class polynomial has degree
$5(p+1)$ and leading coefficient
\begin{equation}\label{eq:momentleading}
\frac{1}{9^{p+1}(p+1)(p+2)}.
\end{equation}
Thus every moment has a complete finite asymptotic expansion with periodic
coefficients. No exceptional corrections are needed for the moments.
\end{theorem}

\begin{proof}
First prove quasipolynomiality. Fix $n\bmod6$, so $k,\ell$ are constant.
For each $v=0,\ldots,k-1$, the representable auxiliary integers in
$[0,B-1]$ have $L=ku+\ell v$, with $0\leq u\leq U_v$ from
\eqref{eq:U}. In residues $1,2$, add the single extra weight at $B+1$.
Consequently, for any $j\geq0$,
\begin{equation}\label{eq:redsums}
\sum_{L\in\mathcal L_n}w(L)^j
=\sum_{v=0}^{k-1}\sum_{u=0}^{U_v}(Au+Cv)^j
 +\eta\,w(B+1)^j.
\end{equation}
When $\eta=0$ the last term is omitted. Every $U_v$, $d$, and $e$ is an
affine function of $n$ on the fixed residue class. All their relevant
values and all displayed ranges are valid for the allowed $n\geq2$.
Expanding powers and summing ordinary powers proves that every power sum
of the full representatives $de\,w+ia+jc$ is a polynomial in $n$ on that
residue class. Formula \eqref{eq:moment} therefore expresses $M_p(n)$ as a
\emph{rational function} of $n$ with rational coefficients.

There is a useful elementary integrality argument to remove the possible
denominator. If $R(q)\in\Q(q)$ is integral for all sufficiently large
integers $q$, perform polynomial division $R=Q+U$, with
$Q\in\Q[q]$ and $U(q)\to0$. A fixed integer $D>0$ clears all coefficients
of $Q$, so $U(q)\in D^{-1}\Z$ at those integer arguments. For large $q$,
$|U(q)|<1/D$, hence $U(q)=0$. A rational function with infinitely many
zeros is identically zero. Thus $R$ is a polynomial. Apply this to
$n=6q+r$ and $M_p(6q+r)$. Since the rational formula was already valid
for every allowed $n\geq2$, its polynomial simplification is valid at
all those indices, not merely eventually.

The degree and leading coefficient follow from the independently proved
limit law in Section~\ref{sec:limit}, which uses only the Ap\'ery normal
form and the Frobenius formula, not the present quasipolynomiality claim.
Specifically it gives
\[
M_p(n)\sim\frac{F_n^{p+1}}{(p+1)(p+2)}
\sim\frac{n^{5p+5}}{9^{p+1}(p+1)(p+2)}.
\]
A polynomial with this positive asymptotic leading term has exactly the
stated degree and leading coefficient in every residue class.
\end{proof}

For example, the first gap-power sum starts
\begin{equation}\label{eq:firstmomexp}
M_1(n)=\frac{n^{10}}{486}+\frac{29n^9}{972}
       +\frac{395n^8}{1944}+c_{n\bmod6}n^7+O(n^6),
\end{equation}
where
\[
(c_0,c_1,c_2,c_3,c_4,c_5)
=\frac1{3888}(3455,3185,3215,3185,3455,3305).
\]
Appendix~\ref{app:firstmoment} gives all terms exactly.
The program \texttt{derive\_moments.py} implements \eqref{eq:redsums}, the
rectangular lift, and \eqref{eq:moment}; it generates the polynomial
for a requested order by symbolic summation, not interpolation. Exact
order-one and order-two outputs accompany the report.

\section{A general gap-limit principle and the triangular law}\label{sec:limit}

\subsection{Transferring a limit of Ap\'ery sets}

Let $S_j$ be a sequence of numerical semigroups with Frobenius numbers
$F_j\to\infty$. Choose any $m_j\in S_j\setminus\{0\}$ such that
$m_j/F_j\to0$. Define the normalized empirical Ap\'ery measure
\[
\nu_j=\frac1{m_j}\sum_{w\in\Ap(S_j;m_j)}\delta_{w/F_j}.
\]
Its support is contained in $[0,1+m_j/F_j]$, since the maximum Ap\'ery
element is $F_j+m_j$.

\begin{samepage}
\begin{theorem}[Gap-limit transfer principle]\label{thm:transfer}
Suppose $\nu_j$ converges weakly to a probability measure $\nu$ on $[0,1]$.
Put $\mu=\int t\dd\nu(t)$. Then $\mu\geq1/2$, and the cumulative
distribution of a uniformly chosen gap divided by $F_j$ converges uniformly
on $[0,1]$ to
\begin{equation}\label{eq:transfer}
G(x)=\frac{\displaystyle\int_0^1\min(x,t)\dd\nu(t)}{\mu}.
\end{equation}
The limit is absolutely continuous and has a nonincreasing density
\begin{equation}\label{eq:transferdensity}
g(x)=\frac{\nu((x,1])}{\mu}\quad\text{for almost every }x\in(0,1),
\qquad 0\leq g(x)\leq2.
\end{equation}
Thus its cumulative distribution is concave.
\end{theorem}
\end{samepage}

\begin{proof}
For $0\leq x\leq1$, count the gaps in each residue class using
\eqref{eq:resgaps}. The number in residue $r$ not exceeding $xF_j$ differs
from $\min(xF_j,w_r)/m_j$ by at most an absolute constant. Hence, uniformly
in $x$,
\begin{equation}\label{eq:count}
\#\{t\in\Gap(S_j):t\leq xF_j\}
=F_j\int\min(x,t)\dd\nu_j(t)+O(m_j).
\end{equation}
At $x=1$ this, or the exact genus formula, gives
$g(S_j)/F_j\to\mu$.
Pairing $t$ and $F_j-t$ shows $g(S_j)\geq(F_j+1)/2$, so $\mu\geq1/2$.

The integrands $\min(x,t)$ are uniformly bounded and uniformly Lipschitz
in both variables on the relevant compact interval. Weak convergence,
a finite mesh in $x$, and the Lipschitz bound therefore give uniform
convergence of their integrals. Divide \eqref{eq:count} by the genus to
obtain \eqref{eq:transfer}. Finally,
\[
\int\min(x,t)\dd\nu(t)=\int_0^x\nu((u,1])\dd u.
\]
This identity follows by integrating the indicator $\ind_{\{u<t\}}$ and
interchanging two nonnegative integrals. It proves absolute continuity,
the density formula, monotonicity, and the bound $g\leq1/\mu\leq2$.
\end{proof}

The transfer principle gives a useful necessary shape constraint for
limiting gap distributions whenever a semigroup modulus is negligible
relative to the Frobenius number. By compactness, every subsequential
limit can be treated after extracting a subsequence of the Ap\'ery
measures. This concerns the \emph{sizes of the gaps}; it is not a
statement about distributions of factorization lengths or about random
choice of a numerical semigroup.

\subsection{Uniform Ap\'ery distribution for square-pyramidal generators}

\begin{theorem}[Quantitative triangular gap law]\label{thm:triangular}
Let $X_n$ be uniformly distributed on $\Gap(S_n)$. There is an absolute
constant $C$ such that for every $n\geq2$,
\begin{equation}\label{eq:triangle}
\boxed{\ \sup_{0\leq x\leq1}
\left|\Prob\!\left(\frac{X_n}{F_n}\leq x\right)-(2x-x^2)\right|
\leq\frac Cn\ }.
\end{equation}
In particular, $X_n/F_n$ converges to the probability density
$2(1-x)$ on $[0,1]$. For each fixed integer $p\geq0$,
\begin{align}
\E\left[\left(\frac{X_n}{F_n}\right)^p\right]
 &=\frac{2}{(p+1)(p+2)}+O_p(n^{-1}),\label{eq:normalizedmom}\\
M_p(n)&=\frac{n^{5p+5}}{9^{p+1}(p+1)(p+2)}
           +O_p(n^{5p+4}).\label{eq:momasympt}
\end{align}
The limiting normalized mean and variance are $1/3$ and $1/18$.
\end{theorem}

\begin{proof}
Use the original modulus $m=b$, so $m=O(n^3)$ and
$m/F_n=O(n^{-2})$. Write a full representative as
$W=de\,w(L)+ia+jc$. Define the residue-independent polynomial
\begin{equation}\label{eq:Tscale}
T_n:=\frac{deAB}{k}
=\frac{n(n+1)(n+2)(2n+1)(2n+3)}{36}.
\end{equation}
Then \eqref{eq:vuw} gives
\[
W=T_n\frac LB+\frac{5deB}{k}v(L)+ia+jc.
\]
Here $v(L)\leq5$, $d,e=O(n)$, $A,C=O(n^2)$,
$B\asymp n$, and $a,c=O(n^3)$, uniformly over all six residues.
Thus the last three terms are $O(n^4)$, while
$F_n\asymp n^5$ and $T_n/F_n=1+O(n^{-1})$. Since
$0\leq L\leq B+1$, it follows that
\begin{equation}\label{eq:coupling}
\left|\frac W{F_n}-\frac LB\right|=O(n^{-1})
\end{equation}
uniformly over \emph{every} representative in \eqref{eq:fullap}.

Each $L$ occurs with exactly the same multiplicity $de$. The empirical
measure on the numbers $L/B$ differs from the uniform grid
$0,1/B,\ldots,(B-1)/B$ by at most one replacement. Its integral against a
bounded Lipschitz function therefore differs from the integral against
uniform measure on $[0,1]$ by $O(1/B)=O(1/n)$.
For the uniformly bounded $1$-Lipschitz functions $t\mapsto\min(x,t)$,
\eqref{eq:coupling} gives, uniformly in $x\in[0,1]$,
\begin{equation}\label{eq:minlimit}
\frac1b\sum_{W\in\Ap(S_n;b)}\min(x,W/F_n)
=\int_0^1\min(x,t)\dd t+O(n^{-1})
=x-\frac{x^2}{2}+O(n^{-1}).
\end{equation}
Apply \eqref{eq:count}. The genus satisfies
$g_n/F_n=1/2+O(n^{-1})$ (indeed the exact formulas give a stronger error).
Division gives \eqref{eq:triangle} for all sufficiently large $n$.
Enlarging the constant covers the finitely many smaller values.

Integrating against $x^p$ yields
$\int_0^1 2x^p(1-x)\dd x=2/((p+1)(p+2))$.
The uniform cumulative-distribution error gives the same $O_p(n^{-1})$
error for each fixed moment; at $p=0$ the identity is exact. Multiply by
$g_nF_n^p$ and use the leading Frobenius and genus terms to prove
\eqref{eq:momasympt}. The mean and variance follow from the first two
normalized moments.
\end{proof}

The limit theorem is not based on a histogram. It follows from a uniform
comparison of every normalized Ap\'ery representative with a regular grid.
The $O(n^{-1})$ estimate is deliberately robust; determining optimal
second-order corrections is a separate research question.

\section{The rational generating function and the minimal recurrence}\label{sec:gf}

Let $\Theta=x\,\mathrm d/\mathrm dx$ and set
\begin{equation}\label{eq:Hperiod}
U(x)=\frac1{1-x},\qquad
H(x)=\frac{3x}{1-x^2}+\frac{2x^2}{1-x^3},
\end{equation}
so that $H(x)=\sum_{n\geq0}h_nx^n$. The correction polynomial is
\begin{equation}\label{eq:Epoly}
E(x)=10x^2+40x^3+315x^5+1612x^{11}+3135x^{17}+2400x^{23}.
\end{equation}

\begin{theorem}[Exact generating function and minimal denominator]\label{thm:gf}
The ordinary generating function $\Phi(x)=\sum_{n\geq2}F_nx^n$ is
\begin{align}\label{eq:GF}
\Phi(x)={}&\frac1{36}
 (4\Theta^5+32\Theta^4+101\Theta^3+175\Theta^2+102\Theta-36)U(x)\notag\\
&-\frac5{36}(\Theta+1)(\Theta+2)(2\Theta+3)H(x)
 +1+2x+E(x).
\end{align}
In reduced form its denominator is exactly
\begin{equation}\label{eq:D}
\boxed{\ D(x)=(1-x)^6(1+x)^4(1+x+x^2)^4\ }.
\end{equation}
Its numerator has degree $41$. The least eventual quasiperiod of $F_n$ is
six, and its minimal eventual homogeneous constant-coefficient recurrence
has order $18$.
\end{theorem}

\begin{proof}
Multiplication of the $n$th coefficient by $n$ is the action of $\Theta$.
Thus \eqref{eq:GF} follows termwise from \eqref{eq:compactP} and
\eqref{eq:F}. The extended quasipolynomial has coefficients $-1,-2$ at
$n=0,1$, respectively, explaining the terms $1+2x$.

The $U$ term has a pole of order six at $x=1$, with nonzero leading
coefficient. The $H$ term has order at most four there and cannot cancel
it. At $x=-1$, the first summand of $H$ has a nonzero simple pole;
applying the degree-three differential operator in \eqref{eq:GF} raises
it to a pole of order four, with nonzero leading term. The second summand
of $H$ is regular there. At each primitive cube root of unity the second
summand similarly produces a pole of order four, while the first is
regular. There are no other poles, and the added polynomials do not affect
them. This proves that \eqref{eq:D} is the exact reduced denominator.

An eventual quasipolynomial of period $q$ has poles only at $q$th roots
of unity. The poles at $-1$ and the primitive cube roots force $6\mid q$;
period six is already supplied by \eqref{eq:F}. The pole multiplicities
likewise force every eventual recurrence denominator to contain
\eqref{eq:D}; its degree is $6+4+8=18$. Conversely this denominator
provides such a recurrence.

The highest-degree polynomial correction is $2400x^{23}$. Multiplying
it by $D(x)$, whose degree is $18$ and leading coefficient is one,
produces the nonzero coefficient $2400x^{41}$. All other terms have lower
degree after multiplication by $D$, proving the numerator degree.
\end{proof}

Writing $D(x)=\sum_{j=0}^{18}d_jx^j$, the coefficients in increasing order
of $j$ are
\begin{equation}\label{eq:Dcoeff}
\begin{split}
(d_0,\ldots,d_{18})={}&(1,2,-1,-8,-9,6,23,16,-14,-32,\\
 &\hspace{15mm}-14,16,23,6,-9,-8,-1,2,1).
\end{split}
\end{equation}
Therefore the explicit recurrence is
\begin{equation}\label{eq:recurrence}
\boxed{\ \sum_{j=0}^{18}d_jF_{n-j}=0\qquad(n\geq42)\ }.
\end{equation}
At $n=41$ the left side is $2400$, so the starting point is sharp for this
recurrence. The absence of primitive sixth-root poles is worth noting:
period six arises by combining frequencies of periods two and three,
not by requiring all sixth roots. A generic denominator
$(1-x^6)^6$ would therefore be far from minimal.

\section{All-orders algebraic inversion}\label{sec:inverse}

\subsection{The continuous branches}

For $h\in\{0,2,3,5\}$, let $N_h(y)$ be the large positive solution of
$P_h(N_h(y))=y$. Set
\begin{equation}\label{eq:tyR}
t=(9y)^{1/5},\qquad
R_h(u)=1+8u+\frac{101-10h}{4}u^2+\frac{175-45h}{4}u^3
 +\frac{102-65h}{4}u^4-\frac{36+30h}{4}u^5.
\end{equation}
Then $9P_h(N)=N^5R_h(1/N)$.

\begin{theorem}[Convergent Puiseux inverse]\label{thm:inverse}
For sufficiently large positive $y$,
\begin{equation}\label{eq:inverseall}
N_h(y)=t-\frac85+\sum_{m\geq1}c_m(h)t^{-m},
\qquad
c_m(h)=-\frac1m[u^{m+1}]R_h(u)^{m/5}.
\end{equation}
This series is convergent, not merely a formal asymptotic expansion.
In particular,
\begin{align}\label{eq:inversefirst}
N_h(y)={}&t-\frac85+\frac{50h+7}{100t}
 -\frac{3(25h+149)}{500t^2}\notag\\
&+\frac{2500h^2-400h+37817}{10000t^3}
 -\frac{3(6250h^2+37125h+17199)}{250000t^4}
 +O(t^{-6}).
\end{align}
There is no $t^{-5}$ term. More generally,
\begin{equation}\label{eq:lacunarity}
\boxed{\ c_{5j}(h)=0\qquad(j\geq1)\ }.
\end{equation}
\end{theorem}

\begin{proof}
Put $z=1/t$ and $u=1/N$. The defining equation becomes
\begin{equation}\label{eq:implicit}
u=z\,R_h(u)^{1/5},
\end{equation}
using the analytic fifth root equal to one at $u=0$. The implicit-function
theorem gives a unique analytic solution $u(z)=z+O(z^2)$ near zero.
Thus $1/u(z)$ has a convergent Laurent expansion with leading term $z^{-1}$.

For clarity, the coefficient calculation can be seen directly from
Lagrange inversion. Write $\varphi(u)=R_h(u)^{1/5}$.
For $m\geq1$, the change of variable $z=u/\varphi(u)$ in coefficient
extraction gives
\begin{align*}
[z^m]\frac1{u(z)}
&=[u^{m+1}]\varphi(u)^m-[u^m]\varphi(u)^{m-1}\varphi'(u)\\
&=-\frac1m[u^{m+1}]\varphi(u)^m.
\end{align*}
The constant coefficient is $-\varphi'(0)=-8/5$.
This proves \eqref{eq:inverseall}; ordinary binomial expansion yields
\eqref{eq:inversefirst}. For $m=5j$, the expression $R_h(u)^j$ is a
polynomial of degree $5j$, so its coefficient at $u^{5j+1}$ vanishes.
This proves \eqref{eq:lacunarity} and explains the $O(t^{-6})$ remainder
after the fourth inverse-power term.
\end{proof}

The missing coefficients are a general consequence of inverting a
polynomial at infinity: the same argument for a monic degree-$d$ polynomial
shows that the coefficients of $t^{-dj}$ vanish. This is a useful
Lagrange-inversion corollary, not a claim of a new general inversion theorem.
For the present problem it provides an all-orders structural check on
symbolic calculations.

\begin{proposition}[A uniform sufficient convergence region]\label{prop:radius}
The Laurent series in \eqref{eq:inverseall} converges for real $t>100$,
uniformly over the four allowed $h$ values. If it is truncated after
$c_m(h)t^{-m}$, with $m\geq0$ and $m=0$ meaning $t-8/5$, the error is at
most
\begin{equation}\label{eq:tail}
\frac{2\cdot100^{m+2}}{t^{m+1}(1-100/t)}.
\end{equation}
This is a conservative bound, not an estimate of the actual convergence
radius or the actual numerical error.
\end{proposition}

\begin{proof}
On $|u|\leq1/50$, the coefficient bounds in \eqref{eq:tyR} give
$|R_h(u)-1|<0.18$ and $|R'_h(u)|<9.2$ for every allowed $h$.
The principal fifth root is analytic on this disk. For $|z|\leq1/100$,
the map $u\mapsto zR_h(u)^{1/5}$ maps the disk strictly into itself and
has derivative in modulus less than $0.03$. Its iterates from zero
converge uniformly, so the fixed point is analytic in $z$.
The analytic function $V(z)=z/u(z)=R_h(u(z))^{-1/5}$ satisfies $|V(z)|<2$.
Cauchy's coefficient estimate on the circle $|z|=1/100$, followed by a
geometric-series bound, bounds the remainder of $z^{-1}V(z)$ after its
$z^m$ term by \eqref{eq:tail}.
\end{proof}

No logarithmic or exponentially small sectors are needed in this problem.
The forward expansion is a finite exact quasipolynomial, and each inverse
branch is an algebraic function with a convergent Puiseux expansion. This
is stronger information than a purely formal transseries ansatz.

\subsection{The exact discrete inverse}

Define
\[
I(y)=\max\{n\geq2:F_n\leq y\},
\]
whenever the set is nonempty. First note that $P_h(X)$ is strictly increasing
for $X\geq2$. Indeed, its derivative is minimized among the allowed $h$
by $h=5$, and
\[
36P'_5(X)=20X^4+128X^3+153X^2-100X-223>0\qquad(X\geq2).
\]
For the last inequality,
$153X^2-100X-223=(X-2)(153X+206)+189>0$.
The entire sequence $F_n$ is also strictly increasing. The finite range
$2\leq n<24$ is checked directly from \eqref{eq:F}; for $n\geq24$,
Appendix~\ref{app:monotone} gives positive-coefficient polynomial
certificates for all six adjacent differences.

\begin{corollary}[Residue-aware inverse rounding]\label{cor:inverse}
For $y\geq F_{29}=2940909$,
\begin{equation}\label{eq:discreteinverse}
\boxed{\ I(y)=\max_{0\leq r<6}
\left\{r+6\floor{\frac{N_{h_r}(y)-r}{6}}\right\}\ },
\qquad (h_0,\ldots,h_5)=(0,3,2,3,0,5).
\end{equation}
\end{corollary}

\begin{proof}
For this range of $y$, every residue class has an eligible index at least
$24+r$. On those tails the exceptional corrections have disappeared, and
the relevant continuous polynomial is strictly increasing. Its largest
eligible integer in residue $r$ is therefore exactly the lattice floor
in \eqref{eq:discreteinverse}. Taking the largest of the six candidates
proves the result; smaller exceptional indices cannot affect the maximum.
\end{proof}

A single continuous inverse followed by ordinary rounding would in general
ignore the periodic correction. Formula \eqref{eq:discreteinverse} keeps
that information. For numerical use, evaluate an inverse expansion to
suggest a candidate and then certify it using exact comparisons with
\eqref{eq:F}; near a lattice boundary, the exact comparison is essential.
The included \texttt{inverse\_index} routine instead uses monotone binary
search and exact integer arithmetic, so it has no floating-point rounding
ambiguity.

\section{Reproducibility and limits of the computational audit}\label{sec:audit}

The accompanying code uses two genuinely different descriptions. The new
normal form constructs a short reduced Ap\'ery set and lifts it. An
independent Dijkstra implementation finds least representatives as shortest
path distances in the residue graph: a generator gives an edge of that
weight from $r$ to $r+g$ modulo the chosen modulus. It knows nothing about
the quasipolynomial or the maximum-candidate table.

\begin{table}[ht]
\centering
\caption{Executed exact checks. These supplement the infinite proofs.}\label{tab:audit}
\begin{tabularx}{\textwidth}{@{}p{0.25\textwidth}X@{}}
\toprule
Range & Check\\
\midrule
$n=2,\ldots,35$ & All 34 values displayed in the inspected OEIS entry.\\
$n=2,\ldots,1000$ & Reduced Ap\'ery set against independent Dijkstra;
 gcd identities; Frobenius and genus lifting.\\
$n=2,\ldots,60$ & Frobenius and genus from the original three generators,
 with no gcd reduction.\\
$n=2,\ldots,30$ & Every lifted Ap\'ery representative against Dijkstra
 modulo the original $b$.\\
$n=2,\ldots,200$ & Constructive representation checks, largest-gap and
 conductor boundary checks; direct membership comparisons through $n=30$.\\
$n=2,\ldots,12$ & Gap moments of orders $0,1,2,3$ against direct dynamic
 programming enumeration of the gaps.\\
$n=42,\ldots,1000$ & The order-18 recurrence; the boundary residual at
 $n=41$ is checked to be $2400$.\\
$n=2,\ldots,300$ & Exact inverse at the left and right ends of each
 interval $[F_n,F_{n+1})$.\\
Symbolic & Six residue identities for the reductions and genus; positivity
 certificates; denominator and coprimality; the first inverse coefficients.\\
\bottomrule
\end{tabularx}
\end{table}

The random membership samples use a fixed seed, so they are reproducible.
The boundary checks are deterministic. The main verification run and the
symbolic derivations for moment orders one and two all passed. The exact
console outputs are preserved in the package. The programs generate a CSV
with $F_n$, genus, and symmetry defect through $n=1000$.

To reproduce the main audit in a Python environment with SymPy installed,
run
\begin{quote}\ttfamily
python verify.py --symbolic\\
python derive\_moments.py --power 1\\
python derive\_moments.py --power 2
\end{quote}
The first command without \texttt{--symbolic} uses only the Python standard
library. The manuscript compiles with two runs of \texttt{pdflatex}.

\paragraph{What was not established by computation.}
No finite range, however large, proves the eventual formula, the completeness
of the exception list, or the limiting distribution. Those claims rest on
the proofs above. No machine-checked theorem prover was run for this report.
No OEIS edit or GitHub commit was submitted. The mathematical claims are
therefore conventional proved statements in an unrefereed manuscript,
with reproducible computational support and an explicitly limited novelty
assessment.

\section{Further research questions}\label{sec:future}

The following directions concern extensions not proved in this report.
They are presented as research questions, not as conjectures whose truth
has been silently assumed.

\paragraph{1. Higher sums of powers.}
Replace $s_n$ by $s_n^{(r)}=\sum_{j=1}^n j^r$ for a fixed $r\geq3$.
Three consecutive values still have gcd one, because consecutive
differences are coprime powers. Determine the adjacent-gcd structure,
minimal quasiperiods, degrees, and effective cutoffs of their Frobenius
quasipolynomials. The present argument suggests searching first for a
reduced relation with bounded coefficients, but it does not predict a
universal degree without further work.

\paragraph{2. More consecutive generators.}
Determine exact Frobenius formulas for
$\langle s_n,s_{n+1},\ldots,s_{n+r}\rangle$ with fixed $r\geq3$.
Does adding a fourth generator lower the growth exponent, alter only the
leading constant, or create new residue-dependent regimes? A useful
solution should include an explicit normal form rather than only invoke
general eventual quasipolynomiality.

\paragraph{3. A classification of bounded-relation families.}
Characterize polynomial triples which, after gcd reductions, satisfy a
relation $kC-\ell A=HB$ with bounded $k,\ell,H$, together with compatible
residue scalings. Develop an effective criterion that recognizes when
Lemma~\ref{lem:aux} yields the full Ap\'ery set. This would turn the
successful structure here into a systematic proof-generation method.

\paragraph{4. Optimal and second-order gap limits.}
Determine the optimal uniform rate in \eqref{eq:triangle}. More finely,
does a residue-dependent second-order expansion exist for the cumulative
gap count, away from boundary layers and arithmetic discontinuities?
The explicit rectangular lift exposes several scales, of orders
$n^{-1}$ and $n^{-2}$ after normalization, whose interactions should be
analyzed rather than averaged away.

\paragraph{5. Realizable limiting gap densities.}
Theorem~\ref{thm:transfer} forces a nonincreasing density bounded by two
when a semigroup modulus is $o(F)$. Which such densities actually arise
from numerical semigroups? What additional restrictions does addition
closure place on the candidate Ap\'ery limit measure? In particular,
characterize the circumstances giving uniform, triangular, or
piecewise-polynomial gap limits.

\paragraph{6. Minimal recurrences for higher moments.}
For each fixed $p$, determine the exact pole orders of the generating
function of $M_p(n)$, not just the generic bound supplied by its degree and
period. The first three coefficients of $g_n$ and of $M_1(n)$ are
residue independent, so systematic cancellations occur. A theorem giving
the minimal recurrence order as a function of $p$ would sharpen the
all-orders moment result.

\paragraph{7. The rest of the semigroup invariants.}
Use the explicit Ap\'ery order to determine pseudo-Frobenius numbers,
semigroup type, and minimal defining relations for all six residue
classes. The symmetry classification is only the first structural
invariant. It remains to identify how the small maximum-switching
exceptions interact with these finer algebraic invariants.

\paragraph{8. Fast rank and selection among the gaps.}
Membership has a bounded-operation description, but a direct use of
\eqref{eq:resgaps} for gap rank still scans many residues. Derive a compact
floor-sum algorithm for counting gaps below an arbitrary threshold and
for selecting the $j$th gap. A useful target is logarithmic-depth integer
arithmetic after parameter preprocessing, with explicit certificates for
the returned rank or gap.

\paragraph{9. Certified inverse evaluation.}
Find the actual nearest complex singularity of each algebraic inverse
branch, replacing the conservative convergence region in
Proposition~\ref{prop:radius}. Combine this with interval arithmetic to
obtain short, guaranteed inverse evaluations even when $y$ is close to a
residue-lattice boundary. The exact polynomial comparisons provide an
independent correctness check.

\paragraph{10. Formalization and an OEIS proof artifact.}
Formalize the rectangular lifting lemma and the least auxiliary residue
principle once, and specialize them to the six rows by exact arithmetic.
A desirable deliverable is a checked theorem giving both membership
certificates and the largest-gap formula, not only a finite table.
Only after independent review should the exact formula and its references
be proposed as an OEIS update. The separate
\texttt{oeis\_submission.txt} file is a draft for that purpose, not a
record of a submitted change.

\appendix
\section{The exact sum of the gaps}\label{app:firstmoment}

For every $n\geq2$, the first gap moment is given by the following four
polynomials. As with the genus, no exceptional correction is necessary.
For $n\equiv0,4\pmod6$,
\begin{align}\label{eq:M10}
M_1(n)=\frac{n(n+1)(n+2)(2n+1)}{3888}\bigl(&4n^6+44n^5+227n^4
 +775n^3\notag\\&+1848n^2+2853n+2322\bigr).
\end{align}
For $n\equiv1,3\pmod6$,
\begin{align}\label{eq:M11}
M_1(n)=\frac{(n+1)(n+2)(n+3)}{3888}\bigl(&8n^7+68n^6+294n^5
 +625n^4\notag\\&+841n^3+441n^2-576n+405\bigr).
\end{align}
For $n\equiv2\pmod6$,
\begin{align}\label{eq:M12}
M_1(n)=\frac{(n+1)(n+2)(2n+5)}{3888}\bigl(&4n^7+36n^6+159n^5
 +371n^4\notag\\&+548n^3+337n^2-474n+96\bigr).
\end{align}
For $n\equiv5\pmod6$,
\begin{align}\label{eq:M15}
M_1(n)=\frac{(n+1)(n+2)}{3888}\bigl(&8n^8+92n^7+498n^6+1627n^5
 +3446n^4\notag\\&+5079n^3+4697n^2+1872n+1425\bigr).
\end{align}

Here is a compact derivation sufficient to reproduce every coefficient.
For $m=b$ and $\mathcal A=\Ap(S_n;b)$, the $p=1$ case of
\eqref{eq:moment} reads
\begin{equation}\label{eq:M1derive}
M_1(n)=\frac1{2b}\sum_{W\in\mathcal A}W^2
        -\frac12\sum_{W\in\mathcal A}W+\frac{b^2-1}{12}.
\end{equation}
Expand $W=de\,w+ia+jc$ and sum over the rectangular ranges. If
$T_j=\sum_{L\in\mathcal L_n}w(L)^j$, then for any $j\geq0$,
\begin{equation}\label{eq:liftsums}
\sum_{W\in\mathcal A}W^j
=\sum_{j_1+j_2+j_3=j}\binom{j}{j_1,j_2,j_3}
 (de)^{j_1}T_{j_1}a^{j_2}c^{j_3}
 \left(\sum_{i=0}^{e-1}i^{j_2}\right)
 \left(\sum_{v=0}^{d-1}v^{j_3}\right).
\end{equation}
Use \eqref{eq:redsums} for $T_0,T_1,T_2$, with
$\sum_{u=0}^q u=q(q+1)/2$ and
$\sum_{u=0}^q u^2=q(q+1)(2q+1)/6$.
Substitution into \eqref{eq:M1derive} yields
\eqref{eq:M10}--\eqref{eq:M15}. The accompanying symbolic script performs
exactly these operations and checks the results against independently
enumerated gap sums. The same multinomial identity is used for every
higher moment.

\section{Positive-coefficient monotonicity certificates}\label{app:monotone}

For $r=0,\ldots,5$, put $n=24+r+6q$ with $q\geq0$.
The exact difference $F_{n+1}-F_n$ is
\[
720q^4+A_rq^3+B_rq^2+C_rq+D_r,
\]
where Table~\ref{tab:monotone} gives positive integer coefficients.
Thus all these differences are positive for $q\geq0$.
The remaining consecutive differences, from $n=2$ through $n=23$,
are positive by direct evaluation of \eqref{eq:F}.

\begin{table}[ht]
\centering
\caption{Exact monotonicity certificates on the six tails.}\label{tab:monotone}
\begin{tabular}{rrrrr}
\toprule
$r$ & $A_r$ & $B_r$ & $C_r$ & $D_r$\\
\midrule
0 & 12348 & 79366 & 226585 & 242435\\
1 & 13068 & 88870 & 268396 & 303739\\
2 & 13428 & 93844 & 291280 & 338798\\
3 & 14148 & 104188 & 340807 & 417830\\
4 & 14148 & 104152 & 340422 & 416808\\
5 & 15228 & 120646 & 424390 & 559300\\
\bottomrule
\end{tabular}
\end{table}

\section{Provenance and package contents}\label{app:package}

The package contains the manuscript source and compiled PDF, the main exact
verification program and its executed output, a symbolic moment generator,
exact first- and second-moment polynomial files and their outputs, a CSV of
values through $n=1000$, a machine-readable numerator and denominator for
$\Phi(x)$, a draft OEIS contribution, and a README with reproduction
instructions. No external paper PDFs or font files are bundled.

The source inspection was performed for this report on October 1, 2026.
The relevant negative finding is limited: the inspected A069762 entry did
not supply the formula, and the cited primary-source search did not locate
it. This report should therefore be cited as an unrefereed explicit
derivation until its proofs and literature status have been independently
reviewed. Its exact assertions do not depend on any empirical extrapolation
or any unproved result in ProveIt.

\begingroup
\raggedright
\begin{thebibliography}{9}

\bibitem{oeis}
The OEIS Foundation Inc.,
\emph{A069762: Frobenius number of the numerical semigroup generated by
three consecutive pyramidal numbers}.
Entry attributed to Victoria A. Sapko (2002), with a displayed extension by
Sean A. Irvine (2024). Inspected October 1, 2026.
\url{https://oeis.org/A069762}.

\bibitem{los}
M. Lepilov, J. O'Rourke, and I. Swanson,
\emph{Frobenius numbers of numerical semigroups generated by three
consecutive squares or cubes},
Semigroup Forum \textbf{91} (2015), 238--259.
DOI: \href{https://doi.org/10.1007/s00233-014-9687-8}{10.1007/s00233-014-9687-8}.
Author manuscript dated October 22, 2014:
\url{https://www.math.purdue.edu/~iswanso/LepilovORourkeSwanson.pdf}.

\bibitem{rr}
A. M. Robles-P\'erez and J. C. Rosales,
\emph{The Frobenius number for sequences of triangular and tetrahedral
numbers}, Journal of Number Theory \textbf{186} (2018), 473--492.
\url{https://arxiv.org/abs/1706.04378}.

\bibitem{rw}
B. H. Roune and K. Woods,
\emph{The parametric Frobenius problem},
Electronic Journal of Combinatorics \textbf{22}(2) (2015), P2.36.
DOI: \href{https://doi.org/10.37236/5112}{10.37236/5112}.
\url{https://arxiv.org/abs/1502.06009}.

\bibitem{shen}
B. Shen,
\emph{The parametric Frobenius problem and parametric exclusion},
arXiv:1510.01349, first posted 2015; version 3 consulted.
\url{https://arxiv.org/abs/1510.01349}.

\bibitem{proveit}
V. Reshetnikov,
\emph{ProveIt: machine-checked mathematics in Lean 4 and Rocq/Coq,
organized by subject}, repository README and Combinatorics directory.
Inspected root-tree snapshot
\texttt{29aca108ed25d714f31b1316512777fbbdc8e006}.
\url{https://github.com/VladimirReshetnikov/ProveIt}.

\end{thebibliography}
\endgroup
\end{document}
